\section{Data and Methodology}
\label{sec:methodology}

The empirical design separates two distinct uses of variance information. The first is an informational allocation channel, in which volatility and Variance Risk Premium variables enter models that determine equity--bond weights. The second is a model-based direct variance-payoff channel, in which lagged implied variance is compared with subsequent realized variance. The two layers are evaluated jointly but are not treated as mechanically comparable investment strategies because they use different payoff mappings and aligned samples.

\subsection{Data, markets, and sample construction}

The analysis covers the United States and Europe. The United States is represented by a broad equity-market series based on the S\&P 500, VIX as the implied-volatility measure, and AGG as the defensive bond asset. Europe is represented by the EURO STOXX 50, VSTOXX/V2TX as the implied-volatility measure, and IEAG.AS as the bond component.

The validated United States pipeline selects SPY for the equity
component, the CBOE VIX (\texttt{\^{}VIX}) for implied volatility,
and AGG for the bond component. The validated European pipeline uses
the EURO STOXX 50 index (\texttt{\^{}STOXX50E}) for equities and
IEAG.AS for bonds. These programmatically downloaded series are
obtained through Yahoo Finance using \texttt{yfinance}; where
applicable, the loader uses automatically adjusted close-like prices.
The European implied-volatility input is a locally harmonized
VSTOXX/V2TX historical series because the required long history is
not reproducible from the current Yahoo Finance interfaces. The raw
local volatility source is therefore treated as an external analytical
input rather than a universally redistributable repository asset.

Because the United States equity component is an adjusted ETF series
whereas the European equity component is an index series, absolute
cross-market return levels should not be interpreted as a perfectly
homogeneous total-return comparison. The principal empirical
inference is conducted within each market against benchmarks
constructed from the same underlying return series.


Raw observations are processed at daily frequency because realized variance requires daily equity returns. Portfolio decisions and performance evaluation are then conducted at monthly frequency. At each month-end, the latest valid volatility features are retained and equity and bond holding-period returns are calculated from month-end prices.

The final validated empirical run ends on 31 July 2026. The United States monthly feature set contains 257 observations and the European feature set contains 195. After the initial estimation requirement and one-step-ahead alignment, the common allocation samples contain 184 observations in the United States and 122 in Europe. The direct variance-payoff panels contain 256 and 194 valid settlement observations, respectively. After lagged risk estimation and common-strategy alignment, the corresponding direct-variance strategy samples contain 232 observations in the United States and 170 in Europe.

The two evidence layers are always evaluated on their own internally aligned samples. Performance statistics from the allocation sample are therefore not mechanically ranked against statistics from the direct-variance sample.

Table~\ref{tab:sample_architecture} summarizes the final sample structure used across the empirical layers.

\begin{table}[htbp]
\centering
\caption{Final empirical sample architecture}
\label{tab:sample_architecture}
\small
\begin{tabular}{lrrll}
\toprule
Evidence layer & US obs. & EU obs. & US start & EU start \\
\midrule
Monthly feature set & 257 & 195 & 2005-03-31 & 2009-05-31 \\
Allocation OOS sample & 184 & 122 & 2011-04-30 & 2015-06-30 \\
Direct payoff diagnostics & 256 & 194 & 2005-04-30 & 2009-06-30 \\
Direct strategy sample & 232 & 170 & 2007-04-30 & 2011-06-30 \\
\bottomrule
\end{tabular}

\vspace{0.4em}
\begin{minipage}{0.92\textwidth}
\footnotesize
All final samples end on 31 July 2026. Allocation and direct-variance statistics are evaluated on their own internally aligned samples and are not mechanically ranked across evidence layers.
\end{minipage}
\end{table}

\subsection{Return construction and variance measures}

Daily equity and bond log returns are calculated from consecutive prices. Let $P_d$ denote the relevant daily price. The log return is

\begin{equation}
r_d
=
\log
\left(
\frac{P_d}{P_{d-1}}
\right).
\end{equation}

Returns spanning more than seven calendar days are set to missing. This prevents a price movement accumulated over a long data interruption from being interpreted as a single-day return. Because realized variance requires a complete 21-observation window, this treatment also prevents long data gaps from contaminating subsequent volatility estimates.

Annualized realized variance is constructed from the most recent 21 valid daily equity log returns:

\begin{equation}
RV_d
=
\frac{252}{21}
\sum_{j=0}^{20}
r_{d-j}^{2}.
\end{equation}

The resulting quantity is a trailing 21-observation realized-variance estimate sampled at month-end. It is not an exact contractual calendar-month variance realization.

Let $V_t$ denote the VIX or VSTOXX index level in volatility percentage points. The implied-variance proxy is

\begin{equation}
IV_t
=
\left(
\frac{V_t}{100}
\right)^2.
\end{equation}

This conversion treats the volatility index as an annualized implied-volatility measure. It does not reconstruct an exact variance-swap strike from the underlying option surface.

Two principal transformations of the Variance Risk Premium are used. The raw proxy is

\begin{equation}
VRP_t
=
IV_t
-
RV_t.
\end{equation}

while the relative transformation is

\begin{equation}
LogVRP_t
=
\log
\left(
\frac{IV_t}{RV_t}
\right).
\end{equation}

Log realized variance and log implied variance are also retained as features. Small positive numerical floors are applied before logarithmic transformation to avoid undefined values.

For allocation models, these variables are observed at month-end and used to determine the following period's portfolio. They are therefore informational state variables rather than derivative returns.

\subsection{Benchmark portfolios and implementation costs}

Three conventional portfolios provide the main allocation benchmarks: buy-and-hold equity, a 60/40 equity--bond portfolio, and an equal-weighted equity--bond portfolio. The static balanced portfolios are rebalanced monthly to their target weights.

The 60/40 benchmark holds 60\% equity and 40\% bonds. The equal-weighted benchmark holds 50\% in each asset. Buy-and-hold equity remains fully invested in the equity component.

Transaction costs are expressed as a proportional charge on turnover. The baseline assumption is 10 basis points per unit of turnover. For the static balanced portfolios, turnover measures the trade required to move from post-return drifted weights back to the target allocation. For the dynamic allocation models, turnover is the sum of the absolute changes in target equity and bond weights.

This benchmark design follows the principle that model complexity should be evaluated relative to transparent and difficult-to-beat alternatives rather than in isolation, consistent with the empirical portfolio-choice evidence discussed in \textcite{DeMiguelGarlappiUppal2009}.

\subsection{Hidden Markov Model allocation}

The first regime framework is a two-state Gaussian Hidden Markov Model (HMM). Let
$S_t \in \{1,2\}$ denote the unobserved regime at month $t$. Regime transitions follow a
first-order Markov chain,

\begin{equation}
\Pr(S_t=j \mid S_{t-1}=i)
=
p_{ij},
\qquad
P
=
\begin{pmatrix}
p_{11} & p_{12} \\
p_{21} & p_{22}
\end{pmatrix},
\qquad
\sum_{j=1}^{2} p_{ij}=1.
\end{equation}

Let $\mathbf{z}_t$ denote the vector of standardized observed features. Conditional on the latent
state, the observation equation is Gaussian,

\begin{equation}
\mathbf{z}_t
\mid
S_t=k
\sim
\mathcal{N}
\left(
\boldsymbol{\mu}_k,
\boldsymbol{\Sigma}_k
\right),
\qquad
k\in\{1,2\},
\end{equation}

where $\boldsymbol{\Sigma}_k$ is restricted to be diagonal. This restriction reduces the number
of covariance parameters and improves stability in repeated out-of-sample estimation.

The parameter vector
$\theta=\{P,\boldsymbol{\mu}_1,\boldsymbol{\mu}_2,
\boldsymbol{\Sigma}_1,\boldsymbol{\Sigma}_2\}$
is estimated by maximum likelihood,

\begin{equation}
\widehat{\theta}
=
\arg\max_{\theta}
\log
p
\left(
\mathbf{z}_1,\ldots,\mathbf{z}_T
\mid
\theta
\right),
\end{equation}

using the expectation--maximization procedure implemented by
\texttt{GaussianHMM}. The implementation uses two components, diagonal covariance matrices,
a maximum of 300 estimation iterations, a convergence tolerance of $10^{-3}$, a minimum
covariance floor of $10^{-5}$, and a fixed random seed for reproducibility.

Three principal information sets are evaluated:

\begin{itemize}
    \item \textbf{HMM RV:}
    monthly equity return and log realized variance;

    \item \textbf{HMM RV + Raw VRP:}
    monthly equity return, log realized variance, and the raw VRP proxy;

    \item \textbf{HMM RV + Log VRP:}
    monthly equity return, log realized variance, and the relative log implied-to-realized
    variance measure.
\end{itemize}

All HMM input variables are standardized inside the estimation sample. Because numerical HMM
state labels have no intrinsic economic meaning, the stress state is identified separately at each
estimation date as the estimated state with the highest average realized variance.

The out-of-sample exercise begins after 72 monthly observations. At each signal date $t$, the
model is re-estimated on an expanding sample containing only observations available through
$t$. The posterior probability of the stress state at the final observation of this sample is therefore
the relevant filtered endpoint probability for the portfolio decision. No observation dated after
$t$ enters the allocation for $t+1$.

Let $p_t^{stress}$ denote this stress probability. The equity allocation is

\begin{equation}
w_t^{eq}
=
0.80
-
0.60
p_t^{stress},
\end{equation}

and the bond allocation is

\begin{equation}
w_t^{bond}
=
1
-
w_t^{eq}.
\end{equation}

Equity exposure therefore varies continuously from 80\% when the estimated stress probability
is zero to 20\% when it is one. The portfolio return is realized during the following month.

\subsection{Markov-switching regression}

The second regime framework is a two-regime Markov-switching regression (RSM). Unlike the
HMM, which models the joint distribution of a feature vector, the RSM directly models monthly
equity returns conditional on the latent regime and a set of exogenous state variables.

Let $S_t\in\{1,2\}$ follow the same first-order Markov transition structure,

\begin{equation}
\Pr(S_t=j\mid S_{t-1}=i)
=
p_{ij}.
\end{equation}

Conditional on regime $S_t=k$, the return equation is

\begin{equation}
R_t^{eq}
=
\alpha_k
+
\boldsymbol{\beta}_k^{\prime}
\mathbf{x}_t
+
\varepsilon_{t,k},
\end{equation}

with

\begin{equation}
\varepsilon_{t,k}
\sim
\mathcal{N}
\left(
0,
\sigma_k^2
\right).
\end{equation}

The intercept $\alpha_k$, all included exogenous coefficients
$\boldsymbol{\beta}_k$, and the conditional variance $\sigma_k^2$ are allowed to switch across
regimes. Monthly equity returns and all exogenous regressors are standardized within each
estimation sample.

Four specifications are estimated:

\begin{itemize}
    \item \textbf{RSM Returns Only:}
    regime-dependent equity-return dynamics without exogenous volatility variables;

    \item \textbf{RSM RV:}
    log realized variance as the exogenous state variable;

    \item \textbf{RSM RV + Raw VRP:}
    log realized variance and the raw VRP proxy;

    \item \textbf{RSM RV + Log VRP:}
    log realized variance and the log implied-to-realized variance ratio.
\end{itemize}

Parameters and transition probabilities are estimated by maximum likelihood using the Markov regression implementation in \texttt{statsmodels}. The numerical procedure uses 10 expectation--maximization
iterations to initialize the likelihood optimization and permits up to 300 subsequent
optimization iterations. Intercepts, exogenous coefficients, and conditional variances are all
specified as regime dependent.

Filtered marginal state probabilities are computed after estimation. As for the HMM, the stress
regime is not assigned from the arbitrary numerical state label: it is defined as the regime with
the highest average realized variance in the corresponding estimation sample.

The RSM uses the same 72-month initial estimation requirement and the same expanding
out-of-sample information set as the HMM. At signal date $t$, estimation uses data only through
$t$, the filtered stress probability at $t$ determines the allocation, and portfolio performance is
realized at $t+1$.

To isolate differences in regime estimation from differences in portfolio construction, the RSM
probability is mapped into the same allocation rule,

\begin{equation}
w_t^{eq}
=
0.80
-
0.60
p_t^{stress},
\qquad
w_t^{bond}
=
1-w_t^{eq}.
\end{equation}

\subsection{Machine-learning stress classification}

The machine-learning layer provides a nonlinear test of whether VRP variables contain
incremental information for predicting adverse market states. Three classifier families are
estimated: L2-penalized logistic regression, Random Forest, and histogram-based Gradient
Boosting. Their role is predictive rather than structural; they do not replace the economic
interpretation of the variance premium. The use of flexible prediction methods is consistent
with the broader asset-pricing evidence discussed by \textcite{GuKellyXiu2020}.

Two feature sets are compared. The \textbf{Base} information set contains monthly equity return,
bond return, log realized variance, and log implied variance. The \textbf{VRP-enhanced} set adds
raw VRP and the log implied-to-realized variance ratio. Each classifier is estimated with both
sets, yielding six ML specifications.

The stress label is generated strictly inside each training sample. An observation at $s+1$ is
classified as stress when either the following-month equity return is at or below the 20th
percentile of the current training distribution or following-month realized variance is at or above
the 80th percentile:

\begin{equation}
Y_{s+1}
=
\mathbf{1}
\left[
R^{eq}_{s+1}
\leq
Q_{0.20}
\ \text{or}\
RV_{s+1}
\geq
Q_{0.80}
\right].
\end{equation}

For the logistic specification, the conditional stress probability is

\begin{equation}
\Pr
\left(
Y_{t+1}=1
\mid
\mathbf{x}_t
\right)
=
\Lambda
\left(
\beta_0
+
\boldsymbol{\beta}^{\prime}
\mathbf{x}_t
\right),
\qquad
\Lambda(a)
=
\frac{1}{1+e^{-a}},
\end{equation}

and the coefficients are estimated with an L2 penalty. The implementation uses
$C=0.50$, balanced class weights, the \texttt{lbfgs} optimizer, and a maximum of 2,000
iterations.

For the Random Forest, the predicted stress probability is the average probability across
$B$ classification trees,

\begin{equation}
\widehat{p}_t^{RF}
=
\frac{1}{B}
\sum_{b=1}^{B}
\widehat{p}_{t,b},
\end{equation}

with $B=200$ trees. Trees are restricted to maximum depth 3, require at least eight
observations per terminal leaf, use square-root feature subsampling, bootstrap samples, and
balanced-subsample class weights. These restrictions deliberately limit tree complexity in the
small rolling samples.

The histogram Gradient Boosting classifier constructs an additive score,

\begin{equation}
F_M(\mathbf{x})
=
F_0(\mathbf{x})
+
\eta
\sum_{m=1}^{M}
f_m(\mathbf{x}),
\end{equation}

which is mapped into a class probability by the logistic link. The implementation uses learning
rate $\eta=0.05$, $M=100$ boosting iterations, maximum tree depth 2, at most seven leaf nodes,
a minimum of ten observations per leaf, L2 regularization equal to 1, balanced class weights,
and no early stopping.

Unlike the HMM and RSM, the ML exercise uses a fixed 72-month rolling training window.
At signal date $t$, features dated $t$ are observable but training feature--label pairs terminate at
$t-1$. Each training label uses the subsequent outcome at $s+1$, and both stress thresholds are
re-estimated using only the current rolling window. The predicted probability at $t$ is therefore
genuinely one-step-ahead.

All three classifier families use the same probability-to-portfolio mapping,

\begin{equation}
w_t^{eq}
=
0.80
-
0.60
\widehat{p}_t^{stress},
\qquad
w_t^{bond}
=
1-w_t^{eq}.
\end{equation}

Prediction quality is evaluated separately through ROC-AUC, precision--recall AUC, Brier
score, log loss, accuracy, balanced accuracy, precision, and recall. Portfolio performance is then
evaluated independently because superior classification does not necessarily imply superior
investor welfare. The final model-family leaders used in the empirical comparison are reported in
Appendix~\ref{app:model-results}.

\subsection{Model-based direct variance-payoff approximation}

The second empirical evidence layer maps variance information more directly into a payoff. It is explicitly a model-based approximation rather than a reconstructed variance-swap backtest.

For settlement month $t$, the variance-strike proxy is lagged implied variance:

\begin{equation}
K_t^{var}
=
IV_{t-1}.
\end{equation}

The settlement proxy is the realized-variance measure observed at month-end $t$:

\begin{equation}
RV_t^{settlement}
=
RV_t.
\end{equation}

The normalized short-variance payoff is then

\begin{equation}
\widetilde{\Pi}_t^{short}
=
\frac{
K_t^{var}
-
RV_t^{settlement}
}{
K_t^{var}
}.
\end{equation}

This quantity is dimensionless. It is not an observed return on invested capital. It is used as the underlying payoff input for the risk-targeted capital mapping.

Four principal specifications are evaluated: a 5\% annual-volatility target that is always active; a 10\% target that is always active; a 10\% target active only when lagged VRP is positive; and a 10\% target active only when lagged VRP exceeds a lagged historical median.

For the positive-VRP specification, exposure at settlement $t$ is permitted when $VRP_{t-1}>0$. The signal is therefore known before $RV_t$ is observed.

The High-VRP rule uses a 36-month historical median with at least 24 observations. The median itself is shifted by one additional month. Consequently, at settlement $t$ the threshold uses entry signals corresponding at the latest to $VRP_{t-2}$. The current entry signal $VRP_{t-1}$ and the current settlement outcome are excluded from threshold estimation.

\subsection{Lagged risk targeting and direct-variance costs}

The payoff-volatility estimate is calculated from the normalized short-variance payoff over a 36-month window, with at least 24 observations, and is shifted by one month. At settlement $t$, the risk estimate therefore uses only payoffs realized through $t-1$.

For annual volatility target $\sigma^{target}$, the monthly target is

\begin{equation}
\sigma^{target}_{m}
=
\frac{
\sigma^{target}
}{
\sqrt{12}
}.
\end{equation}

The unconstrained target notional is

\begin{equation}
N_t^{raw}
=
\frac{
\sigma^{target}_{m}
}{
\widehat{\sigma}_{t-1}
}.
\end{equation}

The implemented notional is non-negative, capped at 0.25, and multiplied by the relevant entry gate. The value 0.25 corresponds to a maximum absolute notional of 25\%.

Each modeled variance exposure represents a one-month position that settles and must be renewed. The implementation cost is therefore charged on the full absolute notional entered each month:

\begin{equation}
Cost_t
=
c
\left|
N_t
\right|.
\end{equation}

The net synthetic direct-variance return is

\begin{equation}
R_t^{DV,net}
=
N_t
\widetilde{\Pi}_t^{short}
-
c
\left|
N_t
\right|,
\end{equation}

where the baseline cost rate $c$ is 10 basis points. The change $|N_t-N_{t-1}|$ is retained only as a diagnostic and is not used as the economic cost base.

\subsection{Performance, welfare, and statistical inference}

Performance is evaluated using annualized geometric return, annualized volatility, Sharpe ratio, Sortino ratio, maximum drawdown, Calmar ratio, empirical 95\% VaR and CVaR, and average turnover.

The Sharpe ratio uses monthly arithmetic mean return, monthly standard deviation, and a zero risk-free rate. The Sortino ratio uses a zero monthly minimum acceptable return. Its downside deviation is calculated over the complete sample, so positive returns contribute zero downside rather than being removed.

Welfare is evaluated using both mean--variance and Constant Relative Risk Aversion certainty equivalents. For monthly return $R_t$ and risk-aversion parameter $\gamma$, the annualized mean--variance certainty equivalent is

\begin{equation}
CEQ^{MV}
=
12
\left(
\overline{R}
-
\frac{\gamma}{2}
s_R^2
\right).
\end{equation}

The analysis considers $\gamma$ values of 1, 3, 5, and 10. CRRA certainty equivalents are calculated from monthly gross returns $1+R_t$ and are annualized geometrically. The CRRA calculation is used only when all gross returns are strictly positive.

For strategy $s$ and benchmark $b$, the welfare difference is

\begin{equation}
\Delta CEQ_{s,b}
=
CEQ_s
-
CEQ_b.
\end{equation}

Fee-equivalent differences are reported in annual basis points by multiplying $\Delta CEQ$ by 10,000.

Sampling uncertainty is assessed using a paired moving-block bootstrap, following the block-resampling logic developed for dependent stationary observations by \textcite{Kunsch1989}. Strategy and benchmark returns are resampled jointly so that cross-strategy dependence is preserved. The baseline procedure uses 2,000 bootstrap replications and six-month blocks. Welfare differences are evaluated against both the 60/40 and equal-weighted benchmarks. A strategy is classified as statistically superior only when the lower bound of the corresponding bootstrap confidence interval is above zero.

\subsection{Robustness design and interpretation}

Robustness is evaluated separately for the allocation and direct-payoff channels. Allocation tests examine transaction-cost sensitivity, no-trade bands, and partial rebalancing. Direct-variance robustness varies transaction costs, the payoff-volatility estimation window, notional caps, subperiods, crisis exclusions, and investor risk aversion.

The direct-variance cost grid contains 0, 10, 25, and 50 basis points. Risk-estimation windows of 12, 24, 36, and 60 months are considered. Maximum notionals of 5\%, 10\%, 15\%, 25\%, and 50\% are tested. Subperiod and crisis-exclusion exercises evaluate whether results are concentrated in a small set of episodes.

Several interpretation rules are maintained throughout the analysis. Allocation and direct-variance strategies are compared only within appropriately aligned samples. The exploratory Pure VRP Proxy is reported separately from the implementable evidence layer. Historical performance is not interpreted as evidence of future profitability.

Most importantly, the direct-variance layer is described throughout as a model-based direct variance-payoff approximation. It does not reconstruct an exact option-replicated variance-swap strike, contractual variance notional, collateral account, variation margin, dealer bid--ask spread, counterparty exposure, or daily mark-to-market path. Its results therefore provide evidence about a stylized variance-carry mechanism rather than a verified historical variance-swap return available to an investor.

\subsection{Code, Data, and Reproducibility}
\label{subsec:reproducibility}

All empirical analyses in this thesis are implemented in Python. The project repository separates data ingestion and feature construction, econometric and machine-learning models, portfolio construction, welfare analysis, robustness testing, reporting, and LaTeX production. The empirical tables and figures reported in the thesis are generated from the same codebase used to estimate the models and construct the portfolios rather than being reconstructed manually for presentation. This emphasis is consistent with the reproducibility concerns documented in empirical finance by \textcite{PerignonEtAl2024}.

The data pipeline combines programmatically downloaded market series with locally supplied files when an external history cannot be reproduced reliably through the same interface. Date parsing, duplicate removal, numerical cleaning, sample filtering, temporal alignment, and calendar-gap treatment are implemented explicitly in the source code. Where redistribution of a raw dataset is restricted or impractical, replication requires access to the original source or an equivalent local input file.

Raw and processed datasets are not assumed to be bundled automatically with the version-controlled repository. Final thesis-facing tables and figures are retained separately so that the numerical results reported in the manuscript can be checked against the outputs produced by the empirical pipeline.

The computational sequence proceeds from baseline market and benchmark construction to HMM and RSM estimation, robustness analysis, welfare analysis, machine-learning stress classification, and the model-based direct variance-payoff extension. Appendix~\ref{app:reproducibility} documents the principal scripts, input locations, output directories, and replication sequence. Appendix~\ref{app:data} documents the principal data-construction diagnostics.
